% !TEX root = ../main.tex
\documentclass[../main.tex]{subfiles}

\begin{document}

\chapter{News Media: The Attributed Map}
\label{app:news-media}

\noindent\textit{News media reveal how others see organizations---their attributed sociotype.
Unlike websites, where organizations curate their own narrative,
news articles reflect how journalists, editors, and sources position actors
in transition discourse.
This appendix documents the bridging workflow that connects news data
to role theory constructs, following the data--theory integration framework
of \textcite{maass2018data}.}

\bigskip


%===============================================================================
\section{The Bridging Problem}
%===============================================================================

News data and role theory address the same phenomenon---who does what
in innovation ecosystems---but from different starting points.
On the data side, NLP pipelines extract entities, verbs, and
semantic relations from news text, producing actor--action--target triples.
On the theory side, role theory specifies how actors are positioned
through attribution: an organization \emph{is} what others \emph{say} it does.
The challenge is bridging between computational extraction
and theoretical attribution---using verb connotation patterns
to operationalize attributed roles, and using role theory
to guide what linguistic features to extract.

\Cref{fig:news-bridging} shows how the \RoleBox{} news pipeline
enacts this bridging workflow.
The upper path (\textit{from patterns to theory}) extracts verb--actor
patterns from news text and maps them onto attributed role constructs.
The lower path (\textit{from theory to data requirements}) uses role
theory to specify what linguistic dimensions and actor features to extract.
The pipeline itself---NER, coreference resolution, dependency parsing,
and RIVETER-X scoring---sits at the intersection as bridging knowledge.

\begin{center}
\begin{tikzpicture}[
    >=stealth,
    cloudstyle/.style={cloud, cloud puffs=12, cloud puff arc=120,
        draw, minimum width=2.8cm, minimum height=1.6cm,
        font=\small\bfseries, text=white, inner sep=0pt},
    blockstyle/.style={rectangle, draw, rounded corners=3pt,
        minimum width=2.6cm, minimum height=1.4cm,
        font=\small\bfseries, text=white, inner sep=4pt},
    bridgestyle/.style={rectangle, draw=AccentCyan, fill=AccentCyan!15,
        rounded corners=3pt, minimum width=3cm, minimum height=1.8cm,
        font=\small\bfseries, text=Slate, inner sep=4pt},
    substyle/.style={font=\tiny, text=Slate, align=left},
    arrowlabel/.style={font=\scriptsize\itshape, text=Slate},
]

% Left: News Data
\node[cloudstyle, fill=Azure!80, draw=Azure] (data) at (-4, 0) {News Data};
\node[substyle, below=0.15cm of data] {
    1{,}273 articles\\
    Entity mentions\\
    Verb relations
};

% Right: Role Theory
\node[blockstyle, fill=AccentPurple!80, draw=AccentPurple] (theory) at (4, 0)
    {Role Theory};
\node[substyle, below=0.15cm of theory] {
    Attributed sociotypes\\
    Power \& agency\\
    Narrative positions
};

% Center: Bridging Pipeline
\node[bridgestyle] (bridge) at (0, 0) {
    \begin{tabular}{c}
    \textsc{RoleBox}\\[-1pt]
    \textsc{News Pipeline}
    \end{tabular}
};

% Upper path: patterns to theory
\draw[->, thick, PandaRose, bend left=20]
    (data.north east) to
    node[arrowlabel, above, pos=0.5] {verb--actor patterns}
    (bridge.north west);
\draw[->, thick, PandaRose, bend left=20]
    (bridge.north east) to
    node[arrowlabel, above, pos=0.5] {attributed roles}
    (theory.north west);

% Lower path: theory to data requirements
\draw[->, thick, CoverGold, bend left=20]
    (theory.south west) to
    node[arrowlabel, below, pos=0.5] {role dimensions}
    (bridge.south east);
\draw[->, thick, CoverGold, bend left=20]
    (bridge.south west) to
    node[arrowlabel, below, pos=0.5] {extraction targets}
    (data.south east);

% Bottom label
\node[font=\small\itshape, text=AccentCyan] at (0, -2.2)
    {Bridging workflow for the attributed map};

\end{tikzpicture}
\captionof{figure}[News media bridging workflow]{%
    Bridging workflow connecting news data to role theory,
    adapted from the data--theory integration framework
    of \textcite{maass2018data}.
    The upper path extracts verb--actor patterns from news text
    and maps them to attributed role constructs;
    the lower path uses theory to guide entity extraction
    and connotation scoring.}
\label{fig:news-bridging}
\end{center}


%===============================================================================
\section{From Patterns to Theory}
\label{sec:news-patterns-theory}
%===============================================================================

The upper path in \cref{fig:news-bridging} traces how raw news text
yields role-relevant patterns.
The workflow proceeds through three stages:
data collection, pattern extraction, and theoretical mapping.

\subsection{Stage 1: Data Collection}

We collected news articles mentioning EIT Knowledge and Innovation
Communities and their partner organizations from GDELT and targeted
web scraping, building a corpus of \textbf{1,273 articles} spanning
2008--2025.\footnote{%
    The initial GDELT queries returned a larger candidate set;
    filtering for relevance, deduplication, and full-text
    availability produced the final 1,273-article corpus
    (\texttt{corpus\_final.json}).}
\Cref{fig:news-collection} summarizes the collection pipeline.

\begin{center}
\begin{tikzpicture}[
    >=stealth,
    box/.style={rectangle, draw=Slate!40, fill=Slate!5,
        rounded corners=2pt, minimum height=0.7cm,
        font=\scriptsize, inner sep=4pt},
    arrow/.style={->, thick, Slate!60},
    statlabel/.style={font=\tiny\itshape, text=Slate!70},
]

% Pipeline stages
\node[box] (seed) at (0,0) {GDELT Queries};
\node[box] (fetch) at (3,0) {Article Fetching};
\node[box] (dedup) at (6,0) {Deduplication};
\node[box] (ner) at (1.5,-1.3) {NER + Coref};
\node[box] (svo) at (4.5,-1.3) {SVO Extraction};
\node[box] (rivet) at (1.5,-2.6) {RIVETER-X Scoring};
\node[box] (agg) at (4.5,-2.6) {Actor Aggregation};

\draw[arrow] (seed) -- (fetch);
\draw[arrow] (fetch) -- (dedup);
\draw[arrow] (dedup) -- (ner);
\draw[arrow] (dedup) -- (svo);
\draw[arrow] (ner) -- (rivet);
\draw[arrow] (svo) -- (rivet);
\draw[arrow] (rivet) -- (agg);

% Stats
\node[statlabel] at (0, -0.4) {partner n-grams};
\node[statlabel] at (6, -0.4) {1{,}273 articles};
\node[statlabel] at (1.5, -1.7) {9{,}269 mentions};
\node[statlabel] at (4.5, -1.7) {2{,}474 relations};

\end{tikzpicture}
\captionof{figure}[News collection pipeline]{%
    Data collection pipeline for news analysis.
    GDELT queries seeded by EIT partner n-grams feed article fetching
    and deduplication, yielding entity mentions and verb relations
    scored on RIVETER-X connotation dimensions.}
\label{fig:news-collection}
\end{center}

The pipeline processes text through six stages:
sentence splitting and cleaning (spaCy),
named entity recognition (spaCy NER + domain gazetteer),
coreference resolution (fastcoref),
dependency-parsed subject--verb--object extraction,
verb connotation scoring (RIVETER-X lexicon),
and cross-document aggregation.\footnote{%
    Coreference resolution is critical for news text.
    A sentence like ``EIT InnoEnergy launched a new accelerator.
    \emph{They} invested 50M euros'' requires linking ``They''
    to EIT InnoEnergy for accurate verb attribution.}

\subsection{Stage 2: Pattern Extraction}

Two families of patterns emerge from the processed corpus.

\paragraph{Verb connotation patterns.}
The RIVETER-X framework scores each verb on six transition-specific
dimensions: Transformation Agency, Innovation Positioning, Legitimacy
Dynamics, Resource Mobilization, Collaboration Capacity, and Systemic
Influence.\footnote{%
    RIVETER-X extends \textcite{antoniak2023riveter}'s verb
    connotation framework with domain-specific dimensions
    calibrated for sustainability transition discourse.}
Aggregating scores by actor reveals systematic differences:
KICs and policy bodies cluster high on legitimacy and resource
mobilization; startups score high on innovation and transformation.

\paragraph{Agent--theme positioning.}
A critical distinction: actors appear as \textbf{agents} (performing
actions) or \textbf{themes} (receiving actions) of verbs.
The agent-to-theme ratio reveals whether news discourse positions
an actor as \emph{acting upon} the world or \emph{being acted upon}.
High agent ratios signal powerful, autonomous positioning
(``X launched,'' ``X transformed'');
high theme ratios signal dependent positioning
(``X was supported by,'' ``X received'').

\subsection{Stage 3: Mapping to Role Constructs}

\Cref{fig:news-mapping} shows how extracted patterns map onto
the theoretical constructs from \cref{ch:cas}.

\begin{center}
\begin{tikzpicture}[
    >=stealth,
    databox/.style={rectangle, draw=Azure!60, fill=Azure!10,
        rounded corners=2pt, font=\scriptsize, inner sep=3pt,
        minimum width=2.2cm, text=Slate},
    theorybox/.style={rectangle, draw=AccentPurple!60, fill=AccentPurple!10,
        rounded corners=2pt, font=\scriptsize, inner sep=3pt,
        minimum width=2.2cm, text=Slate},
    mapstyle/.style={->, thick, PandaRose!70, dashed},
    grouplabel/.style={font=\scriptsize\bfseries, text=Slate},
]

% Data patterns (left)
\node[grouplabel] at (-3, 2.2) {Implicit Patterns};
\node[databox] (d1) at (-3, 1.5) {Verb connotations};
\node[databox] (d2) at (-3, 0.5) {Agent/theme ratios};
\node[databox] (d3) at (-3, -0.5) {Entity co-occurrence};
\node[databox] (d4) at (-3, -1.5) {Temporal frequency};

% Explicit patterns (center)
\node[grouplabel] at (0, 2.2) {Explicit Patterns};
\node[databox, draw=Marigold!60, fill=Marigold!10] (e1) at (0, 1.5)
    {Power/agency profiles};
\node[databox, draw=Marigold!60, fill=Marigold!10] (e2) at (0, 0.5)
    {Active / passive actors};
\node[databox, draw=Marigold!60, fill=Marigold!10] (e3) at (0, -0.5)
    {Discourse coalitions};
\node[databox, draw=Marigold!60, fill=Marigold!10] (e4) at (0, -1.5)
    {Narrative trajectories};

% Theory constructs (right)
\node[grouplabel] at (3, 2.2) {Role Constructs};
\node[theorybox] (t1) at (3, 1.5) {Attributed sociotype};
\node[theorybox] (t2) at (3, 0.5) {Orchestrator / Niche};
\node[theorybox] (t3) at (3, -0.5) {Alliance constellation};
\node[theorybox] (t4) at (3, -1.5) {Role emergence};

% Arrows
\draw[mapstyle] (d1) -- (e1);
\draw[mapstyle] (d2) -- (e2);
\draw[mapstyle] (d3) -- (e3);
\draw[mapstyle] (d4) -- (e4);
\draw[mapstyle] (e1) -- (t1);
\draw[mapstyle] (e2) -- (t2);
\draw[mapstyle] (e3) -- (t3);
\draw[mapstyle] (e4) -- (t4);

% Flow labels
\node[font=\tiny\itshape, text=PandaRose] at (-1.5, -2.2) {explicate};
\node[font=\tiny\itshape, text=PandaRose] at (1.5, -2.2) {embed};

\end{tikzpicture}
\captionof{figure}[Mapping news patterns to role constructs]{%
    From implicit patterns to role constructs, following the
    abstraction pathway of \textcite{maass2018data} Figure~3.
    Implicit patterns (verb connotations, agent/theme ratios) are
    explicated into named role positions, then embedded in
    domain theory.}
\label{fig:news-mapping}
\end{center}

Three mappings deserve emphasis:

\begin{description}[leftmargin=0.5cm, itemsep=0.3em]
\item[Verb connotations $\rightarrow$ Attributed sociotype.]
RIVETER-X dimension profiles distinguish how news positions actors:
KICs score high on legitimacy and resource mobilization
(``fund,'' ``authorize,'' ``support'');
startups score high on transformation and innovation
(``launch,'' ``disrupt,'' ``pioneer'').
The connotation profile is the empirical fingerprint
of the attributed sociotype---what others \emph{say} an actor does.

\item[Agent/theme ratio $\rightarrow$ Structural position.]
Actors with high agent ratios map to orchestrator or
authority positions; those with high theme ratios map to
beneficiary or dependent positions.
This operationalizes the power dimension of role theory:
who acts and who is acted upon.

\item[Co-occurrence networks $\rightarrow$ Discourse coalitions.]
Actors frequently mentioned together in news coverage form
discourse coalitions visible through co-occurrence analysis
\parencite{leifeld2013dna}.
These coalitions map onto the alliance constellations
that CAS theory predicts will emerge around shared functions.
\end{description}


%===============================================================================
\section{From Theory to Data Requirements}
\label{sec:news-theory-data}
%===============================================================================

The lower path in \cref{fig:news-bridging} runs in reverse:
role theory specifies what to look for in news data.
This is the \emph{theory-to-data-requirements} pathway that
\textcite{maass2018data} identify as essential for bridging.

\paragraph{Role attribution theory requires verb-level analysis.}
If roles are attributed through discourse, then the \emph{verbs}
used to describe actors carry the attributional payload.
This requirement guided the adoption of RIVETER-X:
a framework that scores verbs on connotation dimensions
rather than treating all mentions equally.
Not all verbs are equal---``leads'' attributes more agency
than ``participates in.''

\paragraph{Power dynamics theory requires agent--theme parsing.}
The theoretical expectation that intermediaries exercise power
\emph{through} others (enabling, connecting, facilitating) rather
than \emph{over} others required extracting subject--verb--object
triples and distinguishing agent from theme positions.
This guided the dependency parsing and coreference resolution
stages of the pipeline.

\paragraph{Narrative theory requires temporal coverage.}
The Narrative Policy Framework \parencite{jones2014npf} posits
that policy narratives assign hero, villain, and victim roles
that shift over time.
This required building a corpus spanning 2008--2025
to capture how actor portrayals evolve across the EIT's
institutional lifecycle---from establishment through
maturation.


%===============================================================================
\section{The Declared--Attributed Gap}
\label{sec:news-gap}
%===============================================================================

The most theoretically productive pattern is not a data finding
but a \textbf{bridging gap}: the divergence between how organizations
describe themselves on websites (\cref{app:websites}) and how
news media describe them.

This \emph{declared--attributed gap} takes several forms.
Organizations claiming ``disruptive innovation'' on their websites
may be portrayed in news as ``supporting'' or ``participating''---a
downgrade from active to passive positioning.
Conversely, organizations with modest self-descriptions may receive
news coverage attributing them leadership roles they do not claim.

The gap reveals legitimacy dynamics:
successful positioning (claimed roles confirmed by media),
aspirational claiming (claimed roles not yet attributed),
invisible contribution (attributed roles not claimed),
or contested positioning (media attributing roles the actor rejects).

\begin{criticalbox}[The Gap as Bridge to Other Modalities]
The declared--attributed gap cannot be resolved within
website and news data alone.
Resolving it requires triangulation:
Do investment flows confirm the attributed position (\cref{app:crunchbase})?
Does social media discourse reinforce or contest the attribution
(\cref{app:social-media})?
Each resolution path connects this modality's workflow to the next,
creating the multimodal cartography that the dissertation argues for.
\end{criticalbox}


%===============================================================================
\section{Workflow Challenges}
\label{sec:news-challenges}
%===============================================================================

News data presents four categories of methodological challenge,
each requiring specific mitigation strategies.

\paragraph{Curation challenges.}
News corpora are subject to selection bias:
journalists cover ``newsworthy'' actors, leaving routine
activities of peripheral organizations unreported.
Well-resourced PR operations gain disproportionate coverage
\parencite{hager2024media}.
Mitigation: GDELT queries cast a wide net using partner-name
n-grams rather than relying on manually curated outlet lists.

\paragraph{Preprocessing challenges.}
Entity resolution across articles is non-trivial.
The same organization may appear as ``EIT InnoEnergy,''
``InnoEnergy,'' ``the Swedish innovation platform,'' or simply
``the KIC.''
The pipeline combines transformer-based NER with a domain
gazetteer of known EIT organizations and applies coreference
resolution to link pronominal references.

\paragraph{Analysis challenges.}
Verb connotation lexicons do not cover all verbs;
domain adaptation is required for transition-specific
terminology.
The RIVETER-X lexicon extends the base RIVETER lexicon
with 473 unique verbs scored on six dimensions, but gaps
remain for highly technical or neologistic terms.

\paragraph{Archiving challenges.}
News articles are ephemeral: URLs break, paywalls change,
outlets restructure.
The pipeline stores full-text snapshots at collection time
and maintains provenance metadata (URL, date, outlet, fetch
timestamp) to ensure reproducibility even when original
sources become unavailable.


%===============================================================================
\section{Limitations as Bridging Gaps}
\label{sec:news-limitations}
%===============================================================================

News analysis faces six limitations,
each of which motivates a specific cross-modal bridge:

\tableminimal
\begin{center}
\begin{tabular}{@{}p{3cm}p{4cm}p{4cm}@{}}
\toprule
\textbf{Limitation} & \textbf{What it misses} & \textbf{Bridging modality} \\
\midrule
Selection bias & Peripheral actors & Crunchbase (enacted) \\
Source bias & Unpublicized activities & Websites (declared) \\
Framing effects & Ground truth of actions & Investment data \\
Language coverage & Non-English discourse & Wikipedia (multilingual) \\
Temporal lag & Real-time dynamics & Social media \\
Lexicon gaps & Uncovered verb domains & Human-in-the-loop validation \\
\bottomrule
\end{tabular}
\captionof{table}{News limitations and cross-modal bridges}
\label{tab:news-limitations}
\end{center}

These limitations are not reasons to dismiss news data---they are
\emph{bridging gaps} that connect this modality's workflow to the others,
creating the multimodal architecture of \RoleBox{}.


%===============================================================================
\section{Corpus Statistics}
\label{sec:news-corpus-stats}
%===============================================================================

\tableminimal
\begin{center}
\begin{tabular}{@{}lr@{}}
\toprule
\textbf{Metric} & \textbf{Value} \\
\midrule
Articles (corpus\_final) & 1{,}273 \\
Time span & 2008--2025 \\
Entity mentions & 9{,}269 \\
Unique actors (filtered) & 1{,}200+ \\
Verb relations extracted & 2{,}474 \\
Unique verbs & 473 \\
RIVETER-X dimensions & 6 \\
\bottomrule
\end{tabular}
\captionof{table}{EIT news corpus statistics}
\label{tab:news-corpus-stats}
\end{center}

Top extracted actors include EIT (1,457 mentions),
Europe (266), Climate-KIC (128), EIT Health (114),
and EIT Digital (111)---reflecting the corpus focus
on EIT innovation communities.\footnote{%
    The knowledge graph extraction complement
    (kg-gen, 30 articles, 131 entities, 178 relations)
    captures dynamic relationships---investments, launches,
    partnerships---that co-occurrence analysis alone cannot
    reveal. See \texttt{figures/kg\_news\_eit.png}.}


%===============================================================================
\section{Summary}
%===============================================================================

\begin{summarybox}
The news modality produces the \textbf{attributed map}---how others
position organizations through journalistic discourse.
The bridging workflow connects 1,273 news articles to role theory
via two paths: verb connotation patterns map to attributed sociotypes
and power profiles, while agent--theme ratios map to structural
positioning (orchestrator vs.\ beneficiary).
The most productive finding is not within either path but at their
intersection: the \emph{declared--attributed gap} between website
claims and news portrayals that motivates the multimodal
triangulation developed in \cref{app:crunchbase,app:social-media}.
\end{summarybox}

\clearpage
\end{document}
